\documentclass[10pt,a4paper]{amsart}
\usepackage[margin=19mm]{geometry}
\usepackage{amsmath,amssymb,mathtools}
\usepackage[scr=rsfs,cal=euler]{mathalfa}
\usepackage{xfrac}
\usepackage[unicode,hidelinks,bookmarks=false]{hyperref}
\newcommand{\ie}{i.e.\ }
\newcommand{\cf}{cf.\ }
\newcommand{\resp}{respectively\ }
\newcommand{\Subsection}{\subsection}
\providecommand{\nobreakdash}{\nobreak\hbox{-}}
\setlength{\emergencystretch}{2em}
\multlinegap0pt
\usepackage{amssymb}
\usepackage{mathtools}
\usepackage{dsfont}
\usepackage{bm}
\usepackage{float}
\usepackage{xcolor}
\usepackage{tikz}
\newtheorem{prop}{Proposition}
\newtheorem{lemma}{Lemma}
\title{Geometric Approach to Hitchin Components via Higher Complex Structures}
\author{Alexander Thomas}
\date{Source: arXiv:2005.14445v6 (CC BY 4.0)}
\begin{document}
\maketitle

\noindent\emph{Source and licence.} Geometric Approach to Hitchin Components via Higher Complex Structures. Version arXiv:2005.14445v6 (CC BY 4.0), distributed by arXiv under the Creative Commons Attribution 4.0 International licence, http://creativecommons.org/licenses/by/4.0/, as recorded in the arXiv licence metadata for this exact version and shown on the abstract page https://arxiv.org/abs/2005.14445v6. Full licence notice: \texttt{licenses/arxiv-2005.14445-CC-BY-4.0.txt}.\par\smallskip

\noindent\textit{Editorial scope.} Counted unit: the article's prop existence-para-gauge (source0.tex lines 567--569). The article's complete original proof of it is delivered with it (source0.tex lines 571--608, one \texttt{proof} environment of 38 lines). No mathematics is written, repaired or re-derived: the statement and the proof are the article's own lines, in the article's own notation, and no retained line is substituted or re-flowed.  Retained uncounted as the source-local context the proof needs: lines 555--562.  Nothing was omitted from any retained span, so the package carries no omission record.  No retained line contains a citation, so the wrapper carries no bibliography and no bibliography entry.  No retained line contains a figure environment, an image-inclusion command or drawing code, so no image is delivered and the package declares no asset dependency.  Editorial formatting only, in addition: an AMS \texttt{amsart} preamble that restates the article's own declarations and macros used by the retained lines, namely the environments \texttt{prop}, \texttt{lemma} on separate counters, together with the standard packages those lines need.  The retained environments print in this document as Proposition 1, Lemma 1 in order of appearance, whereas the article prints the same items with the numbering of its own full document.  The complete native source of the article as distributed in the arXiv e-print for this exact version is bundled unchanged as \texttt{sources/103769/source0.tex}.
\begin{equation}\label{firstmatrix}
\begin{pmatrix} 
0&  &  & \bm\hat{t}_n\\
1 &0   &  &\vdots \\
& \ddots &\ddots &\bm\hat{t}_2 \\
 &  & 1 & 0 \\
\end{pmatrix},
\end{equation}

\begin{prop}\label{existence-para-gauge}
Let $d+A_1dz+A_2d\bar{z}$ be a connection on $U$. Denote by $a_{i,j}$ the entries of $A_1$. Define the column vector $a=(a_{i,1})_{i=2,...,n}$ and the differential operator $\widetilde{\nabla}=\partial+(a_{i,j})_{2\leq i,j\leq n}$. Then $\det(B)\neq 0$ on $U$ if and only if $\det(a,\widetilde{\nabla}a,...,\widetilde{\nabla}^{n-2}a)\neq 0$.
\end{prop}

\begin{proof}
Since $A_1$ is the matrix of the operator $\nabla$, the condition $\det B\neq 0$ is equivalent to finding a companion matrix \eqref{firstmatrix} in the $\mathcal{P}_L$-gauge orbit of $A_1$.
In equations, we look for $P\in\mathcal{P}_L$ such that $P^{-1}A_1P+P^{-1}\partial P=M$ is a companion matrix. This is equivalent to $\nabla P=PM.$

To set some notations, we write $P$ as:
$$P= \begin{pmatrix}
	p_{11} & p_{12} & \cdots &p_{1n} \\
	0 & &  &  \\
	\vdots & C_2& \cdots& C_n \\
	0 &  &  & 
\end{pmatrix},$$
where $C_k$ denotes the $k$\textsuperscript{th} column in the matrix without the first entry (i.e. a column vector of length $n-1$).
We can then compute $PM$:
$$PM=\begin{pmatrix}p_{12}&\cdots &p_{1n}&*\\ &&&* \\ C_2&\cdots &C_n&\vdots \\ &&&*\end{pmatrix},$$
where the stars denote some entries we will not need. This gives $n^2-n$ equations by the first $n-1$ columns.

Our strategy is the following: first we express $p_{ij}$ for $j>1$ in terms of $a_{ij}$ (the ``constants'') and $p_{11}$ (and its derivatives). Second, we get an expression of $p_{11}$ in terms of $a_{ij}$. In this second step, the condition of the proposition appears.

The matrix equation $\nabla P=PM$ gives for all $i,j$ with $j<n$:
\begin{equation}\label{auxmat}
\sum_{k=1}^n a_{ik}p_{kj}+\partial p_{ij} = p_{i,j+1}.
\end{equation} 
This equation expresses $p_{i,j+1}$ in terms of constants and $(p_{k,j})_{1\leq k\leq n}$. Applying this recursively and since $p_{1,j}=0$ for $j>1$, we get the first step.

To achieve our second goal, we prove the following:
\begin{lemma}
Denote by $P_0$ the square-submatrix of $P$ consisting in the columns $C_2, ..., C_n$. We then have $\det(P_0) = p_{11}^{n-1}\det(a,\widetilde{\nabla}a,...,\widetilde{\nabla}^{n-2}a).$
\end{lemma}
\begin{proof}
The first column of the equation $\nabla P=PM$ gives $C_2=p_{11}a$. From the $k$-th column of $\nabla P=PM$ (with $2\leq k\leq n-1$) and the definition of $\widetilde{\nabla}$, we get $$C_{k+1}=\widetilde{\nabla}(C_k)+p_{1k}\, a.$$
Using operations on columns, we can transform $C_{k+1}$ to $p_{11}\widetilde{\nabla}^{k-1}(a)$ for $1\leq k\leq n-1$. The lemma follows.
\end{proof}

Since $1=\det P =p_{11}\det P_0$, we get $$p_{11} = \det(a,\widetilde{\nabla}a,...,\widetilde{\nabla}^{n-2}a)^{-\frac{1}{n}}.$$
Therefore, if the $L$-parabolic gauge exists, then the previous determinant is non-zero.

Conversely, if this determinant is non-zero, we can compute $P$ from $A_1$, giving a transition matrix to a companion form, in which the basis is of the form $B=(s,\nabla s, ..., \nabla^{n-1}s)$. Since $B$ is a basis, we have $\mathrm{det}(B)\neq 0$.
\end{proof}

\end{document}
